%==============================================================================%
% Chapter 3 – HAETAE and Its NTT
%==============================================================================%

\chapter{The \haetae{} Signature Scheme and Its NTT}
\label{ch:haetae_ntt}

\section{Overview of \haetae{}}

\haetae{}~\cite{HAETAE2022} is a module-lattice based digital signature scheme
submitted to the Korean Post-Quantum Cryptography (KPQC) standardisation
process.  It follows the ``Fiat--Shamir with aborts'' paradigm~\cite{Lyubashevsky2012}:
the signing algorithm samples a masking polynomial, computes a commitment via a
hash function, and rejects the sample if the output leaks information about the
secret key.

\subsection{Algebraic Setting}

\haetae{} operates in the polynomial ring
\begin{equation}
  R_q = \Z_q[X]/(X^n + 1), \qquad n = 256,\quad q = 64513.
\label{eq:ring}
\end{equation}
Keys and signatures are vectors in $R_q^k$ for parameter-set-dependent
ranks $k$.  Polynomial multiplication in $R_q$ is the dominant operation in
key generation, signing, and verification.

\subsection{Role of the NTT}

The NTT converts polynomials from the \emph{coefficient domain} to the
\emph{evaluation domain}, where multiplication is pointwise.  The workflow is:
\begin{equation}
  f \cdot g \;\in\; R_q
  \;\longleftrightarrow\;
  \NTT(f) \odot \NTT(g) \;\in\; \F_q^n,
\label{eq:ntt_mul}
\end{equation}
where $\odot$ denotes componentwise multiplication.  A single polynomial
multiplication thus costs two forward NTTs, one componentwise multiplication,
and one inverse NTT, for overall $O(n \log n)$ field operations instead of
$O(n^2)$.

\section{\haetae{} NTT Parameters}
\label{sec:haetae_params}

\begin{table}[htbp]
\centering
\caption{\haetae{} NTT parameters}
\label{tab:params}
\begin{tabular}{lll}
\toprule
Symbol & Value & Description \\
\midrule
$q$    & $64513$ & Modulus ($= 2^8 \cdot 252 + 1 = 252 \cdot 256 + 1$) \\
$n$    & $256$   & Polynomial degree \\
$\zeta_0$ & $426$ & Primitive $512$th root of unity in $\F_q$ \\
$R$    & $2^{32}$ & Montgomery base \\
$\mathsf{QINV}$ & $940508161$ & $q^{-1} \bmod 2^{32}$ \\
$\mathsf{MONT}$ & $14321$ & $R \bmod q = 2^{32} \bmod q$ \\
$\mathsf{MONTSQ}$ & $4214$ & $R^2 \bmod q = 2^{64} \bmod q$ \\
$R^{-1}$ & $50386$ & $R^{-1} \bmod q$ \\
$f$ & $-29720$ & Inverse NTT scaling constant (in Montgomery form) \\
\bottomrule
\end{tabular}
\end{table}

\subsection{Twiddle Factors}

The forward NTT uses 256 table slots in the reference array
\texttt{zetas[256]}.  Slot 0 is a padding value and the forward loop reads
indices $1,\dots,255$.  For $1 \le k < 256$, the concrete table entry is the
signed 32-bit representative of
$\mathsf{MONT} \cdot \zeta_0^{\brev_8(k)} \bmod q$.

\begin{table}[htbp]
\centering
\caption{First sixteen forward twiddle factors \texttt{zetas[1..16]} (signed Montgomery form)}
\label{tab:zetas}
\begin{tabular}{r|rrrrrrrr}
\toprule
$k$ & 1 & 2 & 3 & 4 & 5 & 6 & 7 & 8 \\
\midrule
\texttt{zetas}$[k]$ & 26964 & -16505 & 22229 & 30746 & 20243 & 19064 & -31218 & 9395 \\
\midrule
$k$ & 9 & 10 & 11 & 12 & 13 & 14 & 15 & 16 \\
\midrule
\texttt{zetas}$[k]$ & -30985 & 22859 & -8851 & 32144 & 13744 & 21408 & 17599 & -16039 \\
\bottomrule
\end{tabular}
\end{table}

The C inverse NTT reads the forward table in reverse order as
\texttt{-zetas[--k]}.  The \jasmin{} extraction stores the same inverse
butterfly constants explicitly as \texttt{jzetas\_inv[0..254]} and reserves
\texttt{jzetas\_inv[255]} for the final scaling constant
$f = n^{-1} \cdot \mathsf{MONT}^2 \bmod q = -29720$.

\section{Reference C Implementation}
\label{sec:c_impl}

The reference C implementation in \texttt{ntt.c} provides the algorithmic
ground truth.  It implements the forward NTT (\Cref{alg:ntt}) and the inverse
NTT with an extra final scaling step.

\subsection{Forward NTT}

\begin{lstlisting}[language=C,
  caption={Forward NTT from \texttt{ntt.c}},
  label={lst:c_ntt}]
void ntt(int32_t a[256]) {
  unsigned int len, start, j, k;
  int32_t t, zeta;

  k = 1;
  for (len = 128; len > 0; len >>= 1) {
    for (start = 0; start < 256; start += 2 * len) {
      zeta = zetas[k++];
      for (j = start; j < start + len; j++) {
        t = montgomery_reduce((int64_t)zeta * a[j + len]);
        a[j + len] = a[j] - t;
        a[j]       = a[j] + t;
      }
    }
  }
}
\end{lstlisting}

The loop nest has three levels:
\begin{enumerate}
  \item \textbf{Outer loop}: $\ell \in \{128, 64, 32, 16, 8, 4, 2, 1\}$,
        the current butterfly half-length.  Each iteration is one NTT stage.
  \item \textbf{Middle loop}: iterates over the $256 / (2\ell)$ groups within
        each stage.
  \item \textbf{Inner loop}: performs $\ell$ butterfly operations within each
        group, using the same twiddle factor \texttt{zeta = zetas[k]}.
\end{enumerate}

Key observations:
\begin{itemize}
  \item After stage $s$ (with $\ell = 2^{7-s}$), the array is arranged so
        that each processed segment represents the checked
        \ec{partial\_ntt} sum for the current segment length.  The formal
        invariant is indexed by \((\ell,\mathit{start},\mathit{bsj})\), not
        by a congruence class modulo \(2^\ell\).

  \item The twiddle counter $k$ advances monotonically from 1 to 255,
        consuming all 255 non-trivial twiddle factors.

  \item No modular reduction is applied at the output of butterfly pairs.
        The checked proof tracks stage-dependent signed bit-width bounds:
        the forward refinement theorem assumes a 16-bit input bound and
        proves a 24-bit output bound.
\end{itemize}

\subsection{Inverse NTT}

\begin{lstlisting}[language=C,
  caption={Inverse NTT from \texttt{ntt.c}},
  label={lst:c_invntt}]
void invntt_tomont(int32_t a[256]) {
  unsigned int start, len, j, k;
  int32_t t, zeta;
  const int32_t f = -29720;   /* mont^2 / 256 mod q */

  k = 256;
  for (len = 1; len < 256; len <<= 1) {
    for (start = 0; start < 256; start = j + len) {
      zeta = -zetas[--k];
      for (j = start; j < start + len; j++) {
        t        = a[j];
        a[j]     = t + a[j + len];
        a[j+len] = t - a[j+len];
        a[j+len] = montgomery_reduce((int64_t)zeta * a[j+len]);
      }
    }
  }
  for (j = 0; j < 256; j++)
    a[j] = montgomery_reduce((int64_t)f * a[j]);
}
\end{lstlisting}

The inverse NTT reverses the butterfly ordering: $\ell$ increases from 1 to
128, and the C twiddle counter starts at 256 so that the first butterfly
uses \texttt{-zetas[255]}.  The final loop multiplies each coefficient by
$f = -29720 \equiv n^{-1} \cdot R^2 \pmod{q}$, with the Montgomery reduction
removing one factor of $R$ and leaving each coefficient multiplied by
$n^{-1} \cdot R \pmod{q}$ (i.e., in Montgomery form).

\subsection{Montgomery Reduction}

\begin{lstlisting}[language=C,
  caption={Montgomery reduction from \texttt{reduce.c}},
  label={lst:c_mont}]
int32_t montgomery_reduce(int64_t a) {
  int32_t t;
  t = (int32_t)a * (int32_t)QINV;
  t = (a - (int64_t)t * Q) >> 32;
  return t;
}
\end{lstlisting}

This implements the signed Montgomery reduction.  The cast \texttt{(int32\_t)a}
extracts the low 32 bits of $a$ (interpreted as a signed 32-bit integer), the
product \texttt{t * Q} is computed as a signed 64-bit value, and the right
shift by 32 is an arithmetic right shift on signed integers.

\subsection{Test Harness}

The file \texttt{test\_haetae\_ntt.c} provides a reference test harness that:
\begin{enumerate}
  \item Generates random polynomial inputs within the valid range.
  \item Compares the C forward NTT with the \jasmin{} forward NTT.
  \item Compares the C inverse NTT with the \jasmin{} inverse NTT.
  \item Checks the combined C and \jasmin{} forward--inverse pipelines agree.
\end{enumerate}

These tests provide empirical C/Jasmin agreement.  The round-trip claim used
for interpretation is the mathematical abstract inverse identity together
with the two separate EasyCrypt end-to-end theorems when the inverse theorem's
input-bound precondition is satisfied; it is not inferred from this test
harness alone, and it is not packaged as a standalone checked pipeline theorem.

\section{Structural Comparison with MLKEM NTT}
\label{sec:haetae_vs_mlkem}

\begin{table}[htbp]
\centering
\caption{Comparison of \haetae{} and MLKEM NTT structures}
\label{tab:haetae_mlkem}
\begin{tabular}{lll}
\toprule
Property & \haetae{} & MLKEM \\
\midrule
Modulus $q$ & $64513$ & $3329$ \\
Degree $n$ & $256$ & $256$ \\
NTT type & Full 256-point negacyclic & Incomplete 128-point \\
Bit-reversal exponent & $\brev_8(k)$ & $\brev_7(k)$ \\
Primitive root & $\zeta_0 = 426$ ($512$th) & $\zeta_0 = 17$ ($256$th) \\
Output permutation & Bit-reversed & Bit-reversed \\
Final inv.\ NTT step & Scale by $n^{-1} R \bmod q$ & Scale by $n^{-1} R \bmod q$ \\
\bottomrule
\end{tabular}
\end{table}

The key structural difference is that the \haetae{} NTT uses twiddle factors
$\zeta^{\brev_8(k)}$ where $k$ ranges from 1 to 255, producing a full 256-point
spectral evaluation.  MLKEM uses twiddle factors $\zeta^{\brev_7(2k+1)}$,
producing only a 128-point evaluation (the polynomial is split into even/odd
parts).  This means the \haetae{} NTT result vector contains evaluation at
all 256 primitive $512$th roots of unity, while the MLKEM result vector
contains evaluations at 128 primitive $256$th roots.

This structural difference invalidates any attempt to directly reuse MLKEM NTT
algebra lemmas for the \haetae{} verification.  New algebraic infrastructure
is required, developed in \ec{NTTFullAlgebra.ec} (\Cref{ch:algebra}).
